\begin{definition}[Actual incidence contraction]%
    \label{def:peps_graph_bond_contraction}
    \lean{TNLean.PEPS.graphBondNetwork}
    \leanok
    Let $\Gamma$ be a finite simple graph with an ordered vertex set $\mathcal V$,
    let $X$ be a finite virtual alphabet, and let $P_v$ be the physical alphabet
    at vertex $v$. For local coefficients $a_v:X^{E(v)}\times P_v\to\mathbb C$,
    define the actual contracted coefficient
    \begin{align}
        \Psi_a(\sigma)=\sum_{\eta:E(\Gamma)\to X}
            \prod_{v\in\mathcal V}a_v(\eta|_{E(v)},\sigma_v).
    \end{align}
    Physical alphabets may depend on the vertex; no auxiliary tensor structure
    or assumed coefficient identity is needed.
\end{definition}

\begin{theorem}[Uniformly numbered tensor coefficients]%
    \label{thm:peps_graph_bond_contraction_tensor}
    \lean{TNLean.PEPS.stateCoeff_groupBondTensor_eq_graphBondNetwork}
    \leanok
    \uses{def:peps_graph_bond_contraction}
    When all physical alphabets are $\{0,\ldots,d-1\}$, numbering the virtual
    alphabet identifies this contraction with the original tensor coefficient.
\end{theorem}
\begin{proof}\leanok
    Reindex the virtual summation by the product of the alphabet numberings.
\end{proof}

\begin{definition}[Regrouping physical coordinates by bonds]%
    \label{def:peps_graph_physical_bond_coordinates}
    \lean{TNLean.PEPS.graphSiteBondEndpointEquiv,
        TNLean.PEPS.graphBondRegrouping,
        TNLean.PEPS.graphBondRegroupingIsometry}
    \leanok
    Orient every edge from its smaller endpoint $t(e)$ to its larger endpoint
    $h(e)$. The site coordinates $\sigma_v:E(v)\to X$ are equivalently the bond
    pairs $\beta_e=(\sigma_{h(e)}(e),\sigma_{t(e)}(e))$.
    This bijection induces a linear equivalence of coefficient spaces and its
    corresponding Hilbert-space coordinate isometry.
\end{definition}

\begin{theorem}[Exact endpoint coordinates and overlap preservation]%
    \label{thm:peps_graph_physical_bond_coordinates}
    \lean{TNLean.PEPS.graphSiteBondEndpointEquiv_symm_tail,
        TNLean.PEPS.graphSiteBondEndpointEquiv_symm_head,
        TNLean.PEPS.graphBondRegrouping_apply,
        TNLean.PEPS.graphBondRegrouping_dotProduct}
    \leanok
    \uses{def:peps_graph_physical_bond_coordinates}
    Recovering the site coordinates reads the second endpoint coordinate at
    each tail and the first at each head. Regrouping evaluates the original
    coefficient at these recovered coordinates and preserves every inner product.
\end{theorem}
\begin{proof}\leanok
    Each half-edge occurs once in the endpoint pairs. Reindex the finite inner
    product by this bijection.
\end{proof}

\begin{definition}[Oriented averaging sites and virtual weights]%
    \label{def:peps_graph_weighted_averaging_site}
    \lean{TNLean.PEPS.graphAveragingSite,
        TNLean.PEPS.graphDress,
        TNLean.PEPS.graphDressedAveragingSite,
        TNLean.PEPS.graphThetaWeightedAveragingSite}
    \leanok
    \uses{def:peps_graph_bond_contraction,def:peps_representation_theta}
    Let $G$ be finite and $U$ a matrix representation on $\mathbb C^X$.
    The local averaging site is $|G|^{-1}\sum_g$ of the product of
    $U(g)$ at head legs and $U(g^{-1})^{\mathsf T}$ at tail legs.
    Inserting a matrix $W$ on every virtual leg means summing the original
    site coefficient against the product of its oriented endpoint entries.
    Apply this insertion to the averaging site to obtain the dressed tensor;
    taking $W=\Theta$ gives the fourth-root-weighted tensor.
    These are the finite-graph local contractions underlying
    \cite[Section~7]{Schuch2010PEPS}, with normalized site averages.
\end{definition}

\begin{theorem}[Actual weighted graph coefficients]%
    \label{thm:peps_graph_weighted_averaging_coherent}
    \lean{TNLean.PEPS.graphDress_one,
        TNLean.PEPS.graphDress_averagingSite,
        TNLean.PEPS.graphBondRegrouping_dressedAveragingSite_coherent,
        TNLean.PEPS.graphBondRegrouping_averagingSite_coherent,
        TNLean.PEPS.graphBondRegrouping_thetaWeightedAveragingSite_coherent}
    \leanok
    \uses{def:peps_graph_weighted_averaging_site,
        thm:peps_graph_physical_bond_coordinates}
    Identity virtual weights leave every site unchanged. If $W$ commutes with
    every $U(g)$, the actual dressed graph state, regrouped by physical bonds, is
    \begin{align}
        \Psi_{U,W}(\beta)=|G|^{-|\mathcal V|}
          \sum_{q:\mathcal V\to G}\prod_{e\in E(\Gamma)}
          [W^2U(q_{h(e)}q_{t(e)}^{-1})]_{\beta_e^1,\beta_e^2}.
    \end{align}
    This includes $W=I$ and the actual fourth-root operator $W=\Theta$.
    No connectedness, valence, or topological assumption is required.
\end{theorem}
\begin{proof}\leanok
    Expand the local convolution to obtain $WU(g^{-1})$ on tail legs and
    $U(g)W$ on head legs. Expand the site averages into one group label per
    vertex and regroup the half-edge product into its endpoint factors.
    The remaining virtual summation factors by edges. Each edge contraction
    gives $U(q_h)W^2U(q_t^{-1})=W^2U(q_hq_t^{-1})$.
    For $W=\Theta$, commutation follows from the character-projector definition.
\end{proof}

\begin{theorem}[Actual finite-graph multiplicity restoration]%
    \label{thm:peps_graph_multiplicity_bond_state}
    \lean{TNLean.PEPS.graphBondRegrouping_blockFourthRootWeight}
    \leanok
    \uses{thm:peps_coherent_multiplicity_transport,
        thm:peps_graph_weighted_averaging_coherent,
        thm:peps_torus_block_multiplicity_state}
    Let $D_i$ be matrix representations of positive dimensions $d_i$, and set
    $U=\bigoplus_iD_i$ and $W=\bigoplus_i d_i^{1/4}I$.
    The actual dressed graph state has matching-sector support on every
    physical bond. The product of the multiplicity-restoring maps
    \begin{align}
        |i;k,l\rangle\longmapsto d_i^{-1/2}
            \sum_{m=1}^{d_i}|i;k,m\rangle\langle i;l,m|
    \end{align}
    sends this state to the actual averaging-site graph state for
    $\bigoplus_i(D_i\otimes I_{d_i})$.
    Support and the state equality are conclusions, not hypotheses.
    This is the finite-simple-graph bond calculation of
    \cite[Section~7]{Schuch2010PEPS}; group-derived Fourier choices and
    full-domain physical isometries are separate consequences.
\end{theorem}
\begin{proof}\leanok
    The two endpoint weights give
    $W^2U(g)=\bigoplus_i\sqrt{d_i}\,D_i(g)$.
    Matching-sector inclusion reconstructs each factor of the actual coherent
    state. The normalized multiplicity map sends this weighted factor to
    $\bigoplus_i(D_i(g)\otimes I_{d_i})$.
    A product physical map acts separately on every factor of every coherent
    summand, giving the asserted actual state.
\end{proof}
